\newpage
\section{微分方程式の解法}
常微分方程式は分類ができ、それぞれ次のような解法がとられる。
\begin{enumerate}[label=\Alph*.]
  \item 非線形
  \begin{enumerate}[label=\arabic*.]
    \item 一階
    \begin{enumerate}[label=\bullet]
      \item 変数分離法
      \item 同次形 \dm{y' = f\Bigl(\frac{y}{x}\Bigr)}
      \item \textbf{ベルヌーイ（Bernoulli）の微分方程式} \dm{p(x)y' + q(x)y = r(x)y^n}
      \item \textbf{完全微分方程式} \dm{P(x,y)dx + Q(x,y)dy = 0}
      \item \textbf{積分因子法} \dm{\grl(x,y)}
      \item \textbf{ラグランジュ（Lagrange）・ダランベール（d'Alembert）の微分方程式} \dm{y = xf(y') + g(y')}
    \end{enumerate}
    \item 二階以上
    \begin{enumerate}[label=\bullet]
      \item 階数低下法（階数低減法）
    \end{enumerate}
  \end{enumerate}
  \vskip1\baselineskip

  \item 線形 \dm{y^{(n)} + p_{n-1}(x)y^{(n-1)} + \cdots + p_{0}(x)y = q(x)}
  \begin{enumerate}[label=\arabic*.]
    \item 一階 \dm{y' + p(x)y = q(x)}
    \begin{enumerate}[label=\roman*.]
      \item 斉次 \dm{y' + p(x)y = 0}
      \begin{enumerate}[label=\bullet]
        \item 変数分離法
      \end{enumerate}
      \item 非斉次
      \begin{enumerate}[label=\bullet]
        \item \textbf{積分因子法} \dm{\grl(x),\,\grl(y)}
        \item \textbf{定数変化法}
      \end{enumerate}
    \end{enumerate}
    \item 二階以上（非定数係数）
    \begin{enumerate}[label=\bullet]
      \item \textbf{オイラー（Euler）の微分方程式} \dm{x^ny^{(n)} + a_{1}x^{n-1}y^{(n-1)} + \cdots + a_{n}y = f(x)}
      % \item スツルム・リウヴィル（Sturm-Liouville）の微分方程式
    \end{enumerate}
    \item 二階以上（定数係数）
    \begin{enumerate}[label=\roman*.]
      \item 斉次
      \begin{enumerate}[label=\bullet]
        \item 特性方程式（補助方程式）
        \item 微分演算子法
      \end{enumerate}
      \item 非斉次
      \begin{enumerate}[label=\bullet]
        \item 特性方程式、未定係数法
        \item 特性方程式、定数変化法
      \end{enumerate}
    \end{enumerate}
  \end{enumerate}
\end{enumerate}

太字部分について以下まとめる。
